\chapter{Szeroka całka Denjoya--Chinczyna}
\providecommand{\atlasref}[2]{#2}
\label{ma1:ch:szeroka-calka-denjoya}
\begingroup
\providecommand{\osc}{\operatorname{osc}}
\providecommand{\Var}{\operatorname{Var}}
\providecommand{\mesh}{\operatorname{mesh}}
\providecommand{\supp}{\operatorname{supp}}
\providecommand{\AC}{\mathrm{AC}}
\providecommand{\BV}{\mathrm{BV}}
\providecommand{\HK}{\mathrm{KH}}
\providecommand{\ind}[1]{\chi_{#1}}
\providecommand{\Dn}{\underline D}
\providecommand{\Du}{\overline D}
\providecommand{\KS}{\mathrm{KS}}
Przyjmujemy oznaczenia i konwencje z rozdziału \ref{ma1:ch:kurzweil-henstock}.
\section{Szeroka całka Denjoya (Denjoya--Chinczyna)}
W porównaniu z wąską całką Denjoya (rozdział
\ref{ma1:ch:waska-calka-denjoya}) zmieniamy dwa elementy
definicji: oscylacje zastępujemy przyrostami między
końcami przedziałów, a zwykłą pochodną --- pochodną
aproksymatywną. Zachowujemy ciągłość funkcji pierwotnej.
Przez otoczenia punktów wnętrza przedziału rozumiemy
otoczenia o promieniu dostatecznie małym, aby mieściły
się w tym przedziale.
\begin{definicja}[Pochodna aproksymatywna]\label{it:def:approx-derivative}
Mierzalny zbiór $E\subset\R$ ma gęstość jeden w punkcie $x$, jeśli $\lambda_1(E\cap(x-r,x+r))/(2r)\to1$ przy $r\downarrow0$. Gęstość zero definiujemy analogicznie, zastępując granicę przez zero. Pochodną aproksymatywną rozpatrujemy w punktach wnętrza przedziału; wartości na końcach nie wpływają na całkę. Dla ciągłej $F:[a,b]\to\R$ i $x\in(a,b)$ mówimy, że $F$ ma skończoną pochodną aproksymatywną $F'_{\mathrm{ap}}(x)=L$, jeżeli dla każdego $\varepsilon>0$ zbiór punktów $y\in[a,b]\setminus\{x\}$ spełniających
$\left|\frac{F(y)-F(x)}{y-x}-L\right|\ge\varepsilon$ ma gęstość zero w $x$. Zbiór ten jest borelowski, ponieważ iloraz różnicowy jest ciągły poza $x$. Zwykła skończona pochodna, o ile istnieje, jest też pochodną aproksymatywną o tej samej wartości: dla dostatecznie małego promienia zbiór złych ilorazów nie przecina otoczenia $x$.
\end{definicja}
\begin{lemat}[Jednoznaczność i liniowość pochodnej aproksymatywnej]\label{it:lem:approx-derivative}
Niech $F,G:[a,b]\to\R$ będą ciągłe i $x\in(a,b)$. Skończona pochodna aproksymatywna w punkcie jest wyznaczona jednoznacznie. Jeśli obie funkcje mają w tym punkcie skończone pochodne aproksymatywne, to ich kombinacja liniowa ma pochodną aproksymatywną równą odpowiedniej kombinacji.
\end{lemat}
\begin{proof}
Gdyby dwie różne liczby $L_1,L_2$ były pochodnymi aproksymatywnymi tej samej funkcji, bierzemy $\varepsilon=|L_1-L_2|/3$. Zbiory punktów, w których iloraz różnicowy różni się od $L_i$ o mniej niż $\varepsilon$, mają gęstość jeden. Ich przecięcie również ma gęstość jeden, ponieważ miara dopełnienia przecięcia jest ograniczona przez sumę miar dopełnień. Przecięcie nie jest więc puste w dostatecznie małych otoczeniach punktu, ale nierówność trójkąta dawałaby $|L_1-L_2|<2\varepsilon$, sprzeczność.
Dla kombinacji $\alpha F+\beta G$ wybieramy dla obu funkcji zbiory dobrych ilorazów z dokładnością $\varepsilon/(1+|\alpha|+|\beta|)$. Na ich przecięciu błąd ilorazu kombinacji jest mniejszy od $\varepsilon$, a przecięcie ma gęstość jeden. To daje tezę.
\end{proof}

\begin{lemat}[Iloczyn, iloraz i złożenie w rachunku aproksymatywnym]
\label{it:lem:approx-derivative-calculus}
Niech $F,G:[a,b]\to\R$ będą ciągłe i niech w punkcie
$x\in(a,b)$ istnieją ich skończone pochodne aproksymatywne.
Wówczas
\[
 (FG)'_{\mathrm{ap}}(x)
 =F(x)G'_{\mathrm{ap}}(x)+G(x)F'_{\mathrm{ap}}(x).
\]
Jeżeli $G(x)\ne0$, to w pewnym otoczeniu $x$ istnieje
iloraz $F/G$ i
\[
 (F/G)'_{\mathrm{ap}}(x)=
 \frac{F'_{\mathrm{ap}}(x)G(x)-F(x)G'_{\mathrm{ap}}(x)}{G(x)^2}.
\]
Jeżeli funkcja $\psi$ jest określona na otoczeniu $F(x)$
i różniczkowalna w $F(x)$, to
\[
 (\psi\circ F)'_{\mathrm{ap}}(x)=\psi'(F(x))F'_{\mathrm{ap}}(x).
\]
W ostatniej tezie zakładamy ponadto ciągłość $\psi$ na
pewnym otoczeniu $F(x)$, aby złożenie było ciągłe lokalnie.
\end{lemat}
\begin{proof}
Piszemy $Q_H(y)=(H(y)-H(x))/(y-x)$ dla $y\ne x$,
$A=F'_{\mathrm{ap}}(x)$ i $B=G'_{\mathrm{ap}}(x)$.
Z tożsamości $Q_{FG}(y)=F(y)Q_G(y)+G(x)Q_F(y)$ mamy
\[
 |Q_{FG}(y)-F(x)B-G(x)A|
 \le |F(y)||Q_G(y)-B|+|G(x)||Q_F(y)-A|
       +|B||F(y)-F(x)|.
\]
Dla $\varepsilon>0$ wybieramy
$\eta=\varepsilon/(2(1+\|F\|_\infty+|G(x)|))$.
Poza sumą dwóch zbiorów gęstości zero pierwsze dwa
składniki łącznie są mniejsze od $\varepsilon/2$.
Ostatni składnik jest mniejszy od $\varepsilon/2$
w dostatecznie małym otoczeniu $x$ dzięki ciągłości $F$.
Dowodzi to reguły iloczynu.

Dla złożenia z definicji różniczkowalności mamy
\[
 Q_{\psi\circ F}(y)=(\psi'(F(x))+e(y))Q_F(y),
 \qquad e(y)\longrightarrow0 \quad(y\to x).
\]
Gdy $F(y)=F(x)$, kładziemy $e(y)=0$; w pozostałych
punktach jest to reszta ilorazu różnicowego $\psi$.
Zbieżność reszty wynika z ciągłości $F$ i różniczkowalności
$\psi$ w $F(x)$. Dla $0<\eta\le1$ poza zbiorem gęstości
zero mamy $|Q_F(y)-A|<\eta$, a stąd
\[
 |Q_{\psi\circ F}(y)-\psi'(F(x))A|
 \le |\psi'(F(x))|\eta+|e(y)|(|A|+1).
\]
Dla ustalonego $\varepsilon>0$ wybieramy
$\eta=\min\{1,\varepsilon/(2(1+|\psi'(F(x))|))\}$,
a następnie promień otoczenia tak, aby
$|e(y)|<\varepsilon/(2(1+|A|))$. W tym otoczeniu
zbiór ilorazów różniących się od żądanej granicy
o co najmniej $\varepsilon$ zawiera się w zbiorze
$\{y\ne x:|Q_F(y)-A|\ge\eta\}$, mającym gęstość zero.
To dowodzi reguły łańcuchowej.
Dla ilorazu ciągłość $G$ i $G(x)\ne0$ zapewniają
niezerowość $G$ w otoczeniu $x$. Stosujemy regułę
złożenia do $\psi(t)=1/t$ i $G$, otrzymując
$(1/G)'_{\mathrm{ap}}(x)=-B/G(x)^2$.
Reguła iloczynu zastosowana lokalnie do $F$ i $1/G$
daje wzór ilorazowy.
\end{proof}

\begin{przyklad}[Pochodna aproksymatywna bez zwykłej pochodnej]
\label{it:ex:approx-not-ordinary}
Na $[-1,1]$ połóżmy $c_n=2^{-n}$, $r_n=2^{-3n-2}$
dla $n\ge1$. Wybieramy funkcję
\[
 \psi(t)=\begin{cases}
 (1-4t^2)^2,& |t|\le1/2,\\
 0,& |t|>1/2.
 \end{cases}
\]
W punktach $\pm1/2$ zarówno wielomian, jak i jego
pochodna $-16t(1-4t^2)$ przyjmują wartość zero.
Funkcja należy więc do $C_c^1((-1,1))$,
spełnia $0\le\psi\le1$ i $\psi(0)=1$.
Przedłużamy $\psi$ przez zero i definiujemy
\[
 F(x)=\sum_{n\ge1}c_n\psi((x-c_n)/r_n).
\]
Przedziały nośników $J_n=[c_n-r_n,c_n+r_n]$
są parami rozłączne: $r_n\le c_n/16$, zatem
$c_{n+1}+r_{n+1}\le17c_n/32<15c_n/16\le c_n-r_n$.
Szereg jest lokalnie skończony poza $0$, gdyż
$c_n+r_n\to0$. Każdy składnik jest ciągły.
W zerze wszystkie składniki są równe zeru.
Dla $\varepsilon>0$ wybieramy $N$ z $c_N<\varepsilon$.
Skończenie wiele przedziałów $J_n$ dla $n<N$
ma dodatnią odległość od zera; wybieramy otoczenie
zera nieprzecinające żadnego z nich (dla $N=1$
nie nakładamy dodatkowego warunku).
Ponieważ przedziały $J_n$ są parami rozłączne,
w tym otoczeniu co najwyżej jeden składnik jest
niezerowy, a jego moduł nie przekracza
$c_N<\varepsilon$. Stąd ciągłość $F$ także w zerze.
Ponieważ $F(c_n)/c_n=1$, a $F(x)/x=0$ dla $x<0$,
zwykła pochodna w $0$ nie istnieje.

Dla $0<\rho\le1/2$ wybieramy $N$ z
$2^{-N-1}<\rho\le2^{-N}$. Gdy $n<N$, mamy
$c_n-r_n\ge15c_n/16\ge(15/8)2^{-N}>\rho$.
Zatem zbiór, na którym $F$ jest niezerowa w
$(-\rho,\rho)$, ma miarę nie większą niż
\[
 \sum_{n\ge N}2r_n
 =\frac{2^{-3N-1}}{1-2^{-3}}=\frac47\,2^{-3N}
 <\frac{32}{7}\rho^3.
\]
Po podzieleniu przez $2\rho$ dostajemy wielkość
dążącą do zera. Dla każdego $\varepsilon>0$ zbiór
punktów z $|F(y)/y|\ge\varepsilon$ zawiera się
w zbiorze, na którym $F$ jest niezerowa. Dlatego
$F'_{\mathrm{ap}}(0)=0$. Przykład rozróżnia oba
pojęcia pochodnej nawet w klasie funkcji ciągłych.
\end{przyklad}

\begin{definicja}[Klasa $\AC G$ i szeroka całka Denjoya]\label{it:def:denjoy-broad}
Dla ciągłej $F:[a,b]\to\R$ i $E\subset[a,b]$ piszemy $F\in\AC(E)$, jeśli dla każdego $\varepsilon>0$ istnieje $\eta>0$ takie, że dla każdej skończonej rodziny przedziałów $[u_j,v_j]$ o rozłącznych wnętrzach i końcach należących do $E$ warunek $\sum_j(v_j-u_j)<\eta$ pociąga $\sum_j|F(v_j)-F(u_j)|<\varepsilon$. Piszemy $F\in\AC G[a,b]$, gdy $[a,b]=\bigcup_{n\ge1}E_n$ dla zbiorów domkniętych $E_n$ z $F\in\AC(E_n)$. Funkcja $f:[a,b]\to\R$ o skończonych wartościach jest całkowalna w szerokim sensie Denjoya, jeżeli istnieje ciągła $F\in\AC G[a,b]$ z $F'_{\mathrm{ap}}=f$ prawie wszędzie. Definiujemy $(\mathrm D_{\mathrm{sz}})\int_a^b f=F(b)-F(a)$. Jednoznaczność tej definicji uzasadnimy w twierdzeniu \ref{it:thm:acg-uniqueness}. Na przedziale jednopunktowym przyjmujemy całkę zero.
Jeśli $f=\widetilde f$ prawie wszędzie i jedna z tych
funkcji jest całkowalna, to ta sama pierwotna realizuje
całkę drugiej: poza sumą dwóch zbiorów wyjątkowych jej
pochodna aproksymatywna jest równa obu funkcjom.
W szczególności zmiana skończonych wartości na zbiorze
miary zero nie zmienia całkowalności ani wartości całki.
\end{definicja}
\begin{twierdzenie}[Domknięcie zbioru i działania algebraiczne w klasie $ACG$]
\label{ma1:kh:thm:acg-algebra}\label{it:thm:acg-closure-algebra}
Dla ciągłej $F:[a,b]\to\R$ i dowolnego $E\subset[a,b]$
zachodzi $F\in\AC(E)$ wtedy i tylko wtedy, gdy $F\in\AC(\overline E)$.
W definicji $\AC G$ można zatem używać dowolnych przeliczalnych
pokryć. Klasa $\AC G[a,b]$ jest algebrą rzeczywistą z jedynką,
zamkniętą także na składanie z zewnętrzną funkcją lipschitzowską
na zakresie $F$.
\end{twierdzenie}
\begin{proof}
Kierunek od $\overline E$ do $E$ wynika z inkluzji.
Odwrotnie, dla $\varepsilon>0$ wybieramy $\eta>0$ z warunku
$\AC(E)$ dla błędu $\varepsilon/4$. Skończoną rodzinę
przedziałów o końcach w $\overline E$, rozłącznych wnętrzach
i sumie długości mniejszej od $\eta/2$ porządkujemy od lewej
do prawej, pomijając przedziały jednopunktowe.
Dzielimy ją na dwie grupy według parzystości indeksów.
W każdej grupie różne przedziały oddziela dodatni odstęp,
gdyż pomiędzy nimi leży przedział drugiej grupy o dodatniej
długości. Przybliżamy wszystkie końce punktami z $E$.
Dla dostatecznie dokładnych przybliżeń przedziały w każdej
grupie nadal mają dodatnie długości i rozłączne wnętrza oraz sumę długości
mniejszą od $\eta$. Suma modułów ich przyrostów jest
mniejsza od $\varepsilon/4$. Z ciągłości $F$ przechodzimy
do granicy w obu skończonych sumach, otrzymując łączne
oszacowanie $\varepsilon/2<\varepsilon$. Domknięcia zbiorów
pokrycia uzasadniają równoważność definicji.

Dla $F,G\in\AC G$ bierzemy przecięcia zbiorów z ich
domkniętych pokryć. Na każdym z nich obie funkcje spełniają
warunek $\AC$. Dla $u<v$ mamy
\[
 |F(v)G(v)-F(u)G(u)|
 \le \|F\|_\infty|G(v)-G(u)|+
 \|G\|_\infty|F(v)-F(u)|.
\]
Normy są skończone z ciągłości na przedziale zwartym.
Wybierając wspólny próg długości dla obu funkcji z błędem
$\varepsilon/(2(1+\|F\|_\infty+\|G\|_\infty))$,
otrzymujemy warunek $\AC$ dla iloczynu.
Dla kombinacji liniowej $sF+tG$ na każdym
zbiorze wspólnego pokrycia wybieramy wspólny próg
długości dla obu funkcji z błędem
$\varepsilon/(1+|s|+|t|)$.
Wówczas suma modułów przyrostów kombinacji liniowej
nie przekracza $|s|$ razy sumy modułów przyrostów $F$
plus $|t|$ razy odpowiedniej sumy dla $G$, zatem jest
mniejsza od $\varepsilon$.
Funkcje stałe mają zerowe przyrosty, w szczególności
funkcja stale równa $1$ jest jedynką algebry. Dla lipschitzowskiej $\psi$ ze stałą $L$ mamy
$|\psi(F(v))-\psi(F(u))|\le L|F(v)-F(u)|$,
co daje warunek $\AC$ na tym samym pokryciu: dla $L>0$
wystarczy próg długości dla błędu $\varepsilon/L$, a dla $L=0$
złożenie jest stałe.
Wszystkie otrzymane funkcje są ciągłe.
\end{proof}
\begin{lemat}[Przedłużenie ze zbioru domkniętego]\label{it:lem:ac-affine-extension}
Jeśli $E\subset[a,b]$ jest domknięty, $F:[a,b]\to\R$ jest ciągła i $F\in\AC(E)$, to $F|_E$ ma bezwzględnie ciągłe przedłużenie $u:[a,b]\to\R$. Ponadto dla każdego zbioru borelowskiego $B\subset E$ zachodzi
\[
 \lambda_1^*(F(B))\le\int_B|u'(x)|\,\mathrm d\lambda_1(x).
\]
W szczególności dla $B=E$ obraz jest zwarty, więc można użyć miary zamiast miary zewnętrznej.
\end{lemat}
\begin{proof}
Dla pustego $E$ bierzemy $u=0$, a dla $E=\{c\}$ bierzemy $u=F(c)$. W obu przypadkach przedłużenie jest bezwzględnie ciągłe, a obraz każdego $B\subset E$ ma miarę zero, więc zachodzi również nierówność z tezy. W pozostałych przypadkach oznaczamy luki wewnętrzne przez $(l_k,r_k)$ i przedłużamy $F$ afinicznie na ich domknięciach, a poza $[\min E,\max E]$ stale. Ciągłość tego przedłużenia można sprawdzić tak jak
w lemacie \ref{ma1:kh:lem:acstar-affine-extension}, bez
założenia o oscylacjach. Podajemy szczegóły.
Poza $E$ funkcja $u$ jest lokalnie afiniczna lub stała.
Ustalmy $x\in E$ i $\varepsilon>0$. Z jednostajnej
ciągłości $F$ wybieramy $\rho>0$ tak, aby
$|F(z)-F(x)|<\varepsilon$ dla $|z-x|<2\rho$.
Jeżeli $|y-x|<\rho$ i $y$ leży w luce wewnętrznej
długości mniejszej niż $\rho$, to oba jej końce są
w odległości mniejszej niż $2\rho$ od $x$.
Wartość $u(y)$ jest kombinacją wypukłą wartości $F$
w tych końcach, więc $|u(y)-F(x)|<\varepsilon$.
Luk o długości co najmniej $\rho$ jest skończenie
wiele, gdyż ich wnętrza są rozłączne i suma długości
nie przekracza $b-a$. Dla każdej z tych luk, której
domknięcie nie zawiera $x$, zmniejszamy otoczenie
$x$, aby nie przecinało jej domknięcia. Dla pozostałych
korzystamy z ciągłości interpolacji afinicznej na
domknięciu luki, która w końcu $x$ przyjmuje wartość
$F(x)$. Analogicznie postępujemy z co najwyżej dwiema
lukami brzegowymi, na których przedłużenie jest stałe.
Dla $y\in E$ korzystamy bezpośrednio z ciągłości $F$.
Skończony wybór dodatkowych ograniczeń na promień
otoczenia daje $|u(y)-u(x)|<\varepsilon$ dla wszystkich
dostatecznie bliskich $y$. Zatem $u$ jest ciągła.

Najpierw wykażemy, że $\sum_k|F(r_k)-F(l_k)|<\infty$. Wybieramy $\eta_0$ z warunku $\AC(E)$ dla błędu $1$. Ponieważ suma długości luk nie przekracza $b-a$, po odrzuceniu skończenie wielu luk suma pozostałych długości jest mniejsza od $\eta_0$. Dla każdej skończonej podrodziny ogona warunek $\AC(E)$ ogranicza sumę modułów przyrostów przez $1$. Supremum po podrodzinach daje zbieżność szeregu.
Stąd funkcja
\[
 w(x)=\sum_k\frac{|F(r_k)-F(l_k)|}{r_k-l_k}\chi_{(l_k,r_k)}(x)
\]
należy do $\mathcal L^1$: jest nieujemna, mierzalna, a z twierdzenia \atlasref{ME-08::thm:sum-nonnegative-integrals}{ME-08: twierdzenie o szeregach nieujemnych funkcji} wynika $\int_{[a,b]}w\,\mathrm d\lambda_1=\sum_k|F(r_k)-F(l_k)|<\infty$.

Dla przedziału $[s,t]$ przecinającego $E$ połóżmy
$p=\min(E\cap[s,t])$ i $q=\max(E\cap[s,t])$.
Ekstrema istnieją, ponieważ przecięcie jest niepuste
i zwarte. Jeśli $s<p$, odcinek $[s,p)$ mieści się
w jednej luce zbioru $E$; na nim $u$ jest afiniczna
lub stała, więc $|u(p)-u(s)|=\int_{[s,p]}w\,\mathrm d\lambda_1$.
Dla $s=p$ obie strony tej równości są zerowe.
Analogicznie $|u(t)-u(q)|=\int_{[q,t]}w\,\mathrm d\lambda_1$.
Ponieważ $u(p)=F(p)$ i $u(q)=F(q)$, z nierówności
trójkąta otrzymujemy
\[
 |u(t)-u(s)|\le |F(q)-F(p)|+
 \int_{[s,p]\cup[q,t]}w\,\mathrm d\lambda_1.
\]
Gdy przedział nie przecina $E$, moduł przyrostu jest równy całce $w$ po nim (również w luce brzegowej, gdzie oba są zerowe).
Dla skończonej rodziny przedziałów o rozłącznych wnętrzach środkowe odcinki $[p,q]$ mają końce w $E$, rozłączne wnętrza i sumę długości nie większą od długości całej rodziny. Przy dostatecznie małej tej sumie ich przyrosty są ograniczone przez $\varepsilon/2$ z warunku $\AC(E)$. Wszystkie pozostałe odcinki mają rozłączne wnętrza i łączną długość nie większą od długości rodziny. Bezwzględna ciągłość całki $w$ (lemat \ref{ma1:kh:lem:integral-ac}) daje dla dostatecznie małej długości ich łączny wkład mniejszy od $\varepsilon/2$. Dowodzi to $u\in\AC[a,b]$.

Z twierdzenia \ref{ma1:kh:thm:ac-lebesgue-primitive} wynika $\operatorname{diam}u(J)\le\int_J|u'|\,\mathrm d\lambda_1$ dla każdego przedziału $J\subset[a,b]$ (także względnie otwartego, przez supremum po jego punktach).
Dla $\varepsilon>0$ z lematu \ref{ma1:kh:lem:integral-ac} wybieramy próg $\eta>0$ dla funkcji $|u'|$ i błędu $\varepsilon$. Ponieważ $B$ jest borelowski i ma skończoną miarę, regularność zewnętrzna (\atlasref{ME-04::thm:lebesgue-outer-regular}{ME-04}) daje względnie otwarty $O\supset B$ z $\lambda_1(O\setminus B)<\eta$. Wówczas $\int_{O\setminus B}|u'|\,\mathrm d\lambda_1<\varepsilon$. Rozkładamy $O$ na jego rozłączne składowe przedziałowe $J_k$, względnie otwarte w $[a,b]$. W rozważanym przypadku $E$ ma co najmniej dwa punkty, więc $a<b$; każda niepusta składowa ma dodatnią długość i zawiera liczbę wymierną w swoim wnętrzu, co daje przeliczalność rodziny. Obrazy $u(J_k)$ są przedziałami z ciągłości $u$, a ich długości (także dla przedziałów niedomkniętych) są równe ich średnicom i nie przekraczają całek $|u'|$ po $J_k$. Ponieważ $F(B)=u(B)\subset\bigcup_k u(J_k)$, podaddytywność miary zewnętrznej daje
\[
 \lambda_1^*(F(B))\le\sum_k\int_{J_k}|u'|\,\mathrm d\lambda_1
 =\int_O|u'|\,\mathrm d\lambda_1
 <\int_B|u'|\,\mathrm d\lambda_1+\varepsilon.
\]
Dowolność $\varepsilon$ kończy dowód.
\end{proof}
\begin{twierdzenie}[Pochodna aproksymatywna i jednoznaczność]\label{it:thm:acg-uniqueness}
Każda ciągła funkcja $F\in\AC G[a,b]$ ma skończoną pochodną aproksymatywną prawie wszędzie. Jeśli $F'_{\mathrm{ap}}=0$ prawie wszędzie, to $F$ jest stała. W konsekwencji szeroka całka Denjoya jest dobrze określona i liniowa.
\end{twierdzenie}
\begin{proof}
Niech $[a,b]=\bigcup_n E_n$ będzie domkniętym pokryciem z definicji, a $u_n$ --- bezwzględnie ciągłym przedłużeniem $F|_{E_n}$ z lematu \ref{it:lem:ac-affine-extension}. Z twierdzenia \atlasref{ME-25::cor:lebesgue-density-25}{ME-25: twierdzenie o punktach gęstości} i twierdzenia \ref{ma1:kh:thm:ac-lebesgue-primitive} wynika, że prawie każdy punkt $x\in E_n\cap(a,b)$ jest punktem gęstości jeden zbioru $E_n$ i punktem istnienia skończonej pochodnej $u_n$. Ustalmy taki punkt oraz $\varepsilon>0$. Dla dostatecznie małego $r>0$ ilorazy różnicowe $u_n$ różnią się od $u_n'(x)$ o mniej niż $\varepsilon$ dla $0<|y-x|<r$. Na $E_n$ wartości obu funkcji są równe, więc zbiór złych ilorazów $F$ w tym otoczeniu zawiera się w $(x-r,x+r)\setminus E_n$. Po podzieleniu jego miary przez $2r$ i przejściu z $r\downarrow0$ otrzymujemy zero. Zatem $F'_{\mathrm{ap}}(x)=u_n'(x)$. Wyjątki ze wszystkich $n$ oraz końce $a,b$ tworzą zbiór miary zero, gdyż jest to przeliczalna suma zbiorów zerowych, a zbiory $E_n$ pokrywają $[a,b]$. Pochodna aproksymatywna ma mierzalnego reprezentanta.
Istotnie, przedłużamy $u_n$ stale poza $[a,b]$ i
rozpatrujemy ciąg ciągłych funkcji
$k(u_n(x+1/k)-u_n(x))$. Jego granica górna jest
borelowska jako funkcja o wartościach w rozszerzonej
prostej rzeczywistej; zastępując wartości nieskończone
zerem, otrzymujemy skończoną funkcję borelowską
$v_n$, równą $u_n'$ w punktach wnętrza, w których
istnieje skończona pochodna. Na zbiorze
$E_n\setminus\bigcup_{j<n}E_j$ bierzemy $v_n$. Te zbiory są borelowskie i tworzą rozłączne pokrycie. Otrzymana funkcja zgadza się z pochodną aproksymatywną prawie wszędzie, zatem dzięki zupełności miary Lebesgue'a każda jej wersja o skończonych wartościach na zbiorze wyjątkowym jest mierzalna w sensie Lebesgue'a.

Jeżeli $F'_{\mathrm{ap}}=0$ prawie wszędzie, to $u_n'=0$ prawie wszędzie na $E_n$. Z lematu \ref{it:lem:ac-affine-extension} wynika $\lambda_1(F(E_n))\le\int_{E_n}|u_n'|\,\mathrm d\lambda_1=0$. Obrazy $F(E_n)$ są zwarte, a $F([a,b])=\bigcup_n F(E_n)$, więc z przeliczalnej podaddytywności $\lambda_1(F([a,b]))=0$. Ponieważ ciągły obraz przedziału jest przedziałem, $F([a,b])$ musi być jednym punktem i $F$ jest stała.

Z twierdzenia \ref{ma1:kh:thm:acg-algebra} klasa $\AC G$ jest zamknięta na kombinacje liniowe.
Jeżeli $F_1,F_2$ są pierwotnymi tej samej funkcji $f$,
to $F_1-F_2\in\AC G$, a z lematu \ref{it:lem:approx-derivative}
jej pochodna aproksymatywna jest zerowa poza sumą dwóch zbiorów wyjątkowych.
Wykazana własność daje $F_1-F_2=\mathrm{const}$, więc przyrosty obu pierwotnych są równe.
Jeżeli natomiast $F,G$ są pierwotnymi $f,g$,
to $\alpha F+\beta G$ jest pierwotną $\alpha f+\beta g$
na mocy tego samego lematu (poza wspólnym zbiorem wyjątkowym).
Odejmując wartości w $a$ i $b$, otrzymujemy liniowość całki.
\end{proof}
\begin{wniosek}[Właściwe rozszerzenie klasy KH]
Każda funkcja całkowalna w sensie KH, Perrona lub wąskim sensie Denjoya jest całkowalna w szerokim sensie Denjoya z tą samą wartością całki. Inkluzja jest właściwa.
\end{wniosek}
\begin{proof}
Na mocy twierdzenia \ref{ma1:kh:thm:kh-acgstar-characterization} pierwotna KH należy do $\AC G_*$ i jej zwykła pochodna jest równa funkcji podcałkowej prawie wszędzie. Dla dowolnego $E$ nierówność $|F(v)-F(u)|\le\osc_{[u,v]}F$ daje $\AC G_*\subset\AC G$, a zwykła pochodna jest pochodną aproksymatywną. Ta sama pierwotna realizuje więc całkę szeroką, z tym samym przyrostem. Dla całek Perrona i wąskiej Denjoya stosujemy odpowiednio twierdzenie \ref{it:thm:perron-hk} i wniosek \ref{it:cor:denjoy-hk}. Właściwość wykaże następny przykład.
\end{proof}
\begin{przyklad}[Oscylacje w lukach zbioru Cantora]\label{ma1:kh:ex:denjoy-strict}\label{it:ex:broad-denjoy-strict}
Niech $K\subset[0,1]$ będzie klasycznym zbiorem Cantora. Na etapie $n\ge1$ usuwanych jest $2^{n-1}$ luk $I_{n,k}=(a_{n,k},b_{n,k})$ długości $3^{-n}$. Definiujemy
\[
F(x)=\begin{cases}
2^{-n}\sin^2\!\bigl(\pi3^n(x-a_{n,k})\bigr),&x\in I_{n,k},\\
0,&x\in K.
\end{cases}
\]
Na domknięciu każdej luki wzór określa funkcję klasy $C^1$ zerującą się na końcach. Aby wykazać ciągłość w $x\in K$, ustalamy $\varepsilon>0$ i wybieramy $N$ z $2^{-N}<\varepsilon$. Na lukach poziomu $n\ge N$ mamy $|F|<\varepsilon$. Luk niższych poziomów jest skończenie wiele: te, których domknięcia nie zawierają $x$, można wykluczyć dostatecznie małym otoczeniem $x$, a dla pozostałych stosujemy ciągłość w ich końcach i $F(x)=0$. Na $K$ funkcja jest stała. To dowodzi ciągłości. Domknięcia wszystkich luk oraz $K$ tworzą przeliczalne domknięte pokrycie; na domknięciach luk zachodzi $\AC$, a na $K$ warunek $\AC(K)$ wynika z zerowych przyrostów. Zatem $F\in\AC G[0,1]$.

Pokażemy, że $F\notin\AC G_*[0,1]$. Gdyby $K$ było pokryte przeliczalnie wieloma domkniętymi zbiorami $E_m\cap K$ z $F\in\AC_*(E_m)$, twierdzenie Baire'a zastosowane w niepustej zwartej przestrzeni metrycznej $K$ daje indeks $m$ i niepusty względnie otwarty zbiór $U\subset K$ zawarty w $E_m$. Dla $x\in U$ wybieramy $\rho>0$ z $K\cap(x-\rho,x+\rho)\subset U$. Podstawowe przedziały konstrukcji Cantora zawierające $x$ mają długości $3^{-r}$; wybierając $r$ z $3^{-r}<\rho$, otrzymujemy taki przedział $J$ z $J\cap K\subset E_m$. Wszystkie luki poziomu $n>r$ zawarte w $J$ mają końce w $E_m$. Na poziomie $r$ przedział $J$ jest jednym przedziałem podstawowym; na każdym kolejnym poziomie liczba jego przedziałów podstawowych podwaja się. Zatem na poziomie $n$ usuwamy w $J$ po jednej luce z każdego z $2^{n-r-1}$ przedziałów poziomu $n-1$. Suma długości wynosi $2^{n-r-1}3^{-n}\to0$, a suma oscylacji jest równa $2^{n-r-1}2^{-n}=2^{-r-1}$, co przeczy definicji $\AC_*(E_m)$.

Niech $f=F'$ na dopełnieniu $K$, a na $K$ określmy $f=0$. W lukach mamy
$f(x)=\pi(3/2)^n\sin(2\pi3^n(x-a_{n,k}))$.
Ponieważ $\lambda_1(K)=0$, $F'_{\mathrm{ap}}=f$ prawie wszędzie. Stąd $(\mathrm D_{\mathrm{sz}})\int_0^1 f=F(1)-F(0)=0$. Gdyby $f$ była całkowalna KH, jej pierwotna $G$ należałaby do $\AC G_*$. Różnica $G-F$ miałaby zerową pochodną aproksymatywną prawie wszędzie, więc jako funkcja klasy $\AC G$ na mocy twierdzenia \ref{it:thm:acg-uniqueness} byłaby stała. Wtedy $F$ różniłaby się od $G$ o stałą i należałaby do $\AC G_*$ przez twierdzenie \ref{it:thm:acstar-closure-algebra}, sprzecznie z wykazanym wynikiem. Ponadto na każdej luce poziomu $n$ całka modułu pochodnej wynosi $2\cdot2^{-n}$: funkcja rośnie od zera do $2^{-n}$ i maleje do zera. Na całym poziomie suma wynosi $2^{n-1}\cdot2^{1-n}=1$. Z twierdzenia \atlasref{ME-08::thm:sum-nonnegative-integrals}{ME-08: twierdzenia o szeregach nieujemnych funkcji} zastosowanego do funkcji $|f|\chi_{I_{n,k}}$ oraz z $f=0$ na $K$ otrzymujemy $\int_0^1|f|\,\mathrm d\lambda_1=\sum_{n\ge1}1=\infty$, więc $f\notin\mathcal L^1$.
\end{przyklad}
\begin{uwaga}[Granica teorii]
Szeroka całka Denjoya odtwarza funkcje pierwotne klasy $\AC G$ z ich pochodnych aproksymatywnych, lecz nie dowolne funkcje ciągłe z pochodnych prawie wszędzie. Funkcja Cantora (\atlasref{ME-17::ex:cantor-measure-17}{ME-17: przykład miary Cantora}) ma pochodną aproksymatywną zero prawie wszędzie, ale nie jest klasy $\AC G$, gdyż jest niestała.
\end{uwaga}

% ============================================================
% UZUPELNIENIE MATEMATYCZNE: ROZDZIAL 9
% ============================================================
\subsection*{Różnica między $\AC(E)$ i $\AC_*(E)$}
\begin{przyklad}[Brak kontroli oscylacji na lukach]\label{prz:ac-kontra-acstar}
Niech $K$ będzie klasycznym zbiorem Cantora, a $F$ funkcją z przykładu \ref{ma1:kh:ex:denjoy-strict}. Na $K$ zachodzi $F=0$, więc $F\in\AC(K)$: przyrosty pomiędzy dowolnymi punktami $K$ są równe zeru. Nie zachodzi natomiast $F\in\AC_*(K)$. Istotnie, dla luk $I_{n,k}$ usuwanych na poziomie $n$ mamy końce w $K$ oraz
\[
\sum_{k=1}^{2^{n-1}}|I_{n,k}|=2^{n-1}3^{-n}\longrightarrow0,
\qquad
\sum_{k=1}^{2^{n-1}}\osc_{\overline I_{n,k}}F
=2^{n-1}2^{-n}=\frac12.
\]
Dla każdego $\eta>0$ znajdziemy zatem rozłączne przedziały o końcach w $K$ i łącznej długości mniejszej niż $\eta$, na których suma oscylacji pozostaje równa $1/2$. To pokazuje, dlaczego w szerokiej teorii kontrolujemy jedynie przyrosty, a w wąskiej całe oscylacje.
\end{przyklad}
\begin{twierdzenie}[Wzór Newtona--Leibniza dla szerokiej całki Denjoya]\label{tw:denjoy-szeroka-ftc}
Niech $F\in C[a,b]\cap\AC G[a,b]$. Wówczas $F'_{\mathrm{ap}}$ istnieje i jest skończona prawie wszędzie. Ustalmy dowolną funkcję $f:[a,b]\to\R$ o skończonych wartościach, równą $F'_{\mathrm{ap}}$ wszędzie tam, gdzie ta pochodna istnieje poza pewnym zbiorem miary zero. Wówczas dla dowolnych $a\le u\le v\le b$ funkcja $f$ jest całkowalna w szerokim sensie Denjoya na $[u,v]$ i
\[
 (\mathrm D_{\mathrm{sz}})\int_u^v f(x)\,\mathrm dx=F(v)-F(u).
\]
\end{twierdzenie}
\begin{proof}
Istnienie pochodnej aproksymatywnej prawie wszędzie ustalono w twierdzeniu \ref{it:thm:acg-uniqueness}. Dla $u=v$ obie strony są zerowe, więc dalej zakładamy $u<v$. Jeśli $[a,b]=\bigcup_n E_n$ jest domkniętym pokryciem z $F\in\AC(E_n)$, to zbiory $E_n\cap[u,v]$ są domknięte w $[u,v]$ i pokrywają ten przedział; dlatego $F|_{[u,v]}\in\AC G[u,v]$. W punktach wewnętrznych $[u,v]$ pochodna aproksymatywna ograniczenia zgadza się z pochodną aproksymatywną $F$, o ile istnieje. Definicja szerokiej całki Denjoya daje tożsamość. Jej jednoznaczność wynika z twierdzenia \ref{it:thm:acg-uniqueness}.
\end{proof}
\begin{stwierdzenie}[Lokalność i liniowość szerokiej całki]\label{st:denjoy-szeroka-lokalnosc}
Jeśli $f,g$ są całkowalne w szerokim sensie Denjoya na $[a,b]$ i $\alpha,\beta\in\R$, to $\alpha f+\beta g$ jest całkowalna i jej całka jest odpowiednią kombinacją liniową. Jeśli $a<c<b$, to $f$ jest całkowalna na $[a,b]$ wtedy i tylko wtedy, gdy jest całkowalna na $[a,c]$ i $[c,b]$; wówczas całki są addytywne.
\end{stwierdzenie}
\begin{proof}
Dla pierwotnych $F,G$ bierzemy $H=\alpha F+\beta G$. Przez przecięcia domkniętych zbiorów pokrycia uzyskujemy $H\in\AC G$, a pochodna aproksymatywna jest liniowa wszędzie, gdzie obie pochodne istnieją. Dla lokalności ograniczamy pierwotną do podprzedziałów. Odwrotnie, dwie pierwotne na podprzedziałach można skleić, dodając jednej z nich stałą, aby wartości w $c$ się zgadzały. Dokładniej, dla pierwotnych $F_1$ na $[a,c]$ i $F_2$ na $[c,b]$ kładziemy $F=F_1$ na $[a,c]$ i $F=F_2+F_1(c)-F_2(c)$ na $[c,b]$. Wartości w $c$ są zgodne, więc $F$ jest ciągła. Zbiory z obu domkniętych pokryć pozostają domknięte w $[a,b]$, a przyrosty po dodaniu stałej nie zmieniają się. Połączenie obu rodzin świadczy o przynależności $F$ do $\AC G$. Pochodna aproksymatywna jest równa $f$ poza sumą lokalnych zbiorów wyjątkowych i zbioru $\{a,c,b\}$, ponieważ w punktach wnętrza każdej części zależy tylko od dostatecznie małego otoczenia. Zatem $F$ realizuje całkę na całości, a jej przyrost jest sumą przyrostów po obu częściach.
\end{proof}
\begin{twierdzenie}[Oszacowanie miary obrazu przez pochodną aproksymatywną]
\label{it:thm:acg-image-bound}
Niech $F\in\AC G[a,b]$ i niech $f$ będzie skończonym
mierzalnym reprezentantem $F'_{\mathrm{ap}}$.
Dla każdego zbioru $B\subset[a,b]$ mierzalnego w sensie
Lebesgue'a zachodzi
\[
 \lambda_1^*(F(B))\le \int_B |f|\,\mathrm d\lambda_1,
\]
przy czym prawa strona może być nieskończona.
W szczególności obraz zbioru, na którym $f=0$ prawie
wszędzie, ma miarę zero, niezależnie od miary samego
zbioru.
\end{twierdzenie}
\begin{proof}
Niech $E_n$ będzie domkniętym pokryciem z definicji
$\AC G$, a $u_n$ bezwzględnie ciągłym przedłużeniem
$F|_{E_n}$ z lematu \ref{it:lem:ac-affine-extension}.
Z dowodu twierdzenia \ref{it:thm:acg-uniqueness} mamy
$u_n'=f$ prawie wszędzie na $E_n$.
Wybieramy borelowski zbiór $D\supset B$ taki, że
$\lambda_1(D\setminus B)=0$; można go otrzymać jako przecięcie względnie otwartych
zbiorów $O_k\supset B$ z $\lambda_1(O_k\setminus B)<2^{-k}$,
istniejących na mocy \atlasref{ME-04::thm:lebesgue-outer-regular}{ME-04: regularności zewnętrznej}.
Wtedy $D=\bigcap_k O_k$ jest borelowski i zawiera $B$.
Dla każdego $k$ mamy $D\setminus B\subset O_k\setminus B$,
więc $\lambda_1(D\setminus B)\le2^{-k}$; przechodząc do granicy,
otrzymujemy $\lambda_1(D\setminus B)=0$. Definiujemy rozłączne
zbiory borelowskie
\[
 D_n=D\cap\bigl(E_n\setminus\bigcup_{j<n}E_j\bigr).
\]
Ich suma jest równa $D$. Z lematu
\ref{it:lem:ac-affine-extension}, ponieważ $D_n\subset E_n$,
otrzymujemy
\[
 \lambda_1^*(F(D_n))
 \le\int_{D_n}|u_n'|\,\mathrm d\lambda_1
 =\int_{D_n}|f|\,\mathrm d\lambda_1.
\]
Zatem z monotoniczności i przeliczalnej podaddytywności
miary zewnętrznej oraz przeliczalnej addytywności całki
funkcji nieujemnej wynika
\[
 \lambda_1^*(F(B))
 \le\sum_{n\ge1}\lambda_1^*(F(D_n))
 \le\sum_{n\ge1}\int_{D_n}|f|\,\mathrm d\lambda_1
 =\int_D|f|\,\mathrm d\lambda_1
 =\int_B|f|\,\mathrm d\lambda_1.
\]
Ostatnia równość wynika z zerowej miary $D\setminus B$
oraz addytywności całki funkcji nieujemnej, także gdy całka jest nieskończona.
Wniosek dla zbiorów, na których $f=0$ prawie wszędzie,
jest bezpośredni.
\end{proof}

\begin{wniosek}[Własność Luzina i ograniczone wahanie]\label{it:cor:acg-bv}
Każda funkcja $F\in\AC G[a,b]$ ma własność Luzina (N): obraz zbioru miary Lebesgue'a zero ma miarę zero. Ponadto
\[
 \AC G[a,b]\cap\BV[a,b]=\AC[a,b].
\]
Jeśli $f$ jest całkowalna w szerokim sensie Denjoya, a $F(x)=(\mathrm D_{\mathrm{sz}})\int_a^x f$, to $F\in\BV[a,b]$ wtedy i tylko wtedy, gdy $f\in\mathcal L^1[a,b]$. Wówczas $V_a^b(F)=\int_a^b|f|\,\mathrm d\lambda_1$.
\end{wniosek}
\begin{proof}
Dla zbioru $N$ miary Lebesgue'a zero stosujemy twierdzenie
\ref{it:thm:acg-image-bound} do $B=N$. Całka modułu
mierzalnego reprezentanta po $N$ jest zerowa, więc
$\lambda_1^*(F(N))=0$. Z zupełności miary Lebesgue'a
obraz jest mierzalny i ma miarę zero.

Jeśli $F\in\AC G\cap\BV$, z rozkładu Jordana i twierdzenia o różniczkowaniu funkcji niemalejących (argument z lematu \ref{ma1:kh:lem:bvstar-null-to-acstar}) wynika, że zwykła pochodna $F'$ istnieje prawie wszędzie i jest całkowalna w sensie Lebesgue'a po uzupełnieniu na zbiorze wyjątkowym. Zwykła pochodna jest również aproksymatywna. Funkcja
\[
 G(x)=F(a)+\int_a^x F'(t)\,\mathrm d\lambda_1(t)
\]
jest bezwzględnie ciągła, a więc należy do $\AC G$, i ma zwykłą pochodną $F'$ prawie wszędzie (twierdzenie \ref{ma1:kh:thm:ac-lebesgue-primitive}). Stąd $F-G$ należy do $\AC G$ i ma zerową pochodną aproksymatywną prawie wszędzie. Twierdzenie \ref{it:thm:acg-uniqueness} daje $F-G$ stałą; wartość w $a$ jest zerowa, więc $F=G\in\AC$. Inkluzja odwrotna wynika z definicji oraz ograniczonego wahania funkcji bezwzględnie ciągłej.

Dla pierwotnej szerokiej całki Denjoya mamy $F'_{\mathrm{ap}}=f$ prawie wszędzie. Jeśli $F\in\BV$, pierwsza część daje $F\in\AC$, a zwykła pochodna $F'$ jest równa $f$ prawie wszędzie przez jednoznaczność pochodnej aproksymatywnej (lemat \ref{it:lem:approx-derivative}). Zatem $f\in\mathcal L^1$. Odwrotnie, jeśli $f\in\mathcal L^1$, jego pierwotna Lebesgue'a jest bezwzględnie ciągła i ma pochodną aproksymatywną równą $f$ prawie wszędzie. Jednoznaczność pierwotnych szerokiej całki Denjoya daje jej równość z $F$. Wzór na wahanie wynika teraz z twierdzenia \ref{it:thm:denjoy-bv-ac}, ponieważ $F$ jest bezwzględnie ciągła, a więc jest także pierwotną wąskiej całki Denjoya.
\end{proof}

% Synchronizacja z CAL-09; zachowane etykiety źródłowe.
\begin{twierdzenie}[Porządek i całkowalność bezwzględna w szerokim sensie Denjoya]\label{ma1:kh:thm:wide-denjoy-positive}
Jeśli $f\ge0$ prawie wszędzie i $f$ jest całkowalna w szerokim sensie Denjoya, to $f\in\mathcal L^1[a,b]$. W konsekwencji szeroka całka Denjoya jest monotoniczna, a całkowalność w tym sensie zarówno $f$, jak i $|f|$ jest równoważna $f\in\mathcal L^1[a,b]$.
\end{twierdzenie}
\begin{proof}
Niech $F\in\AC G$ będzie pierwotną $f$. Z twierdzenia
\ref{it:thm:acg-uniqueness} funkcja $f$ jest mierzalna
w sensie Lebesgue'a. Wybieramy mierzalny zbiór zerowy $N$
zawierający końce przedziału i wszystkie wyjątki, tak że
poza $N$ istnieje skończona pochodna aproksymatywna
$F'_{\mathrm{ap}}=f\ge0$. Niech
$Z=\{x\in[a,b]\setminus N:f(x)=0\}$.
Zbiory $N$ i $Z$ są mierzalne. Na mocy twierdzenia
\ref{it:thm:acg-image-bound} zarówno $F(N)$, jak i $F(Z)$
mają miarę zewnętrzną zero: odpowiednie całki $|f|$
są zerowe.

Przypuśćmy, że $F(s)>F(t)$ dla $s<t$. Dla każdego $z\in(F(t),F(s))$ niech $x$ będzie pierwszym punktem $[s,t]$, w którym $F(x)=z$. Istnienie i minimum wynikają z ciągłości oraz zwartości. Mamy $s<x<t$ oraz $F(y)>z$ dla $s\le y<x$ (w przeciwnym razie twierdzenie o wartości pośredniej dałoby wcześniejsze trafienie w $z$). Zatem ilorazy $(F(y)-F(x))/(y-x)$ są ujemne po lewej stronie $x$. Jeśli $x\notin N$, pochodna aproksymatywna nie może być dodatnia, gdyż wtedy dla $\varepsilon=F'_{\mathrm{ap}}(x)/2$ zbiór złych ilorazów zawiera całe dostatecznie małe lewe półotoczenie, o gęstości co najmniej $1/2$. Ponieważ pochodna jest nieujemna poza $N$, musi być zerowa, czyli $x\in Z$. Otrzymujemy
\[
 (F(t),F(s))\subseteq F(N)\cup F(Z),
\]
co przeczy dodatniej długości przedziału, gdyż prawa strona ma miarę zewnętrzną zero. Stąd $F$ jest niemalejąca, zatem ma ograniczone wahanie. Wniosek \ref{it:cor:acg-bv} daje $F\in\AC$ i $f\in\mathcal L^1$. Jeżeli $g\le h$ prawie wszędzie i obie funkcje są całkowalne w szerokim sensie Denjoya, to $h-g$ jest całkowalna na mocy liniowości i należy do $\mathcal L^1$ z pierwszej części. Jej całka jest nieujemna ze zgodności z całką Lebesgue'a, co daje monotoniczność. Jeżeli $f$ i $|f|$ są całkowalne w szerokim sensie Denjoya, to zastosowanie pierwszej części do $|f|$ daje $\int_{[a,b]}|f|\,\mathrm d\lambda_1<\infty$, a $f$ jest mierzalna jako reprezentant pochodnej aproksymatywnej. Odwrotnie, dla $f\in\mathcal L^1$ również $|f|\in\mathcal L^1$ i obie funkcje są całkowalne w szerokim sensie Denjoya dzięki inkluzji klasy Lebesgue'a w KH i KH w szeroką klasę Denjoya.
\end{proof}
\begin{wniosek}[Twierdzenia graniczne dla szerokiej całki Denjoya]\label{ma1:kh:cor:wide-denjoy-limits}
Niech $f_n$ będą całkowalne w szerokim sensie Denjoya i $f_n\to f$ prawie wszędzie, gdzie $f$ ma skończone wartości. Jeśli $f_n\uparrow f$ prawie wszędzie i całki $f_n$ są ograniczone z góry, albo $g\le f_n\le h$ prawie wszędzie dla dwóch funkcji $g,h$ całkowalnych w szerokim sensie Denjoya, to $f$ jest całkowalna w tym sensie, a całki zbiegają do całki $f$. Analogicznie, jeśli $f_n\downarrow f$ prawie wszędzie i całki $f_n$ są ograniczone z dołu, to zachodzi ta sama teza. Ponadto we wszystkich trzech przypadkach
\[
 f_n-f\in\mathcal L^1[a,b],\qquad
 \int_a^b|f_n-f|\,\mathrm d\lambda_1\longrightarrow0.
\]
Jeśli $F_n(x)=(\mathrm D_{\mathrm{sz}})\int_a^x f_n$
oraz $F(x)=(\mathrm D_{\mathrm{sz}})\int_a^x f$, to
\[
 V_a^b(F_n-F)=\int_a^b|f_n-f|\,\mathrm d\lambda_1
 \longrightarrow0.
\]
W szczególności $F_n\to F$ jednostajnie, a całki po
wszystkich podprzedziałach zbiegają jednostajnie względem
ich obu końców.
\end{wniosek}
\begin{proof}
Wszystkie funkcje całkowalne w szerokim sensie Denjoya
są mierzalne na mocy twierdzenia \ref{it:thm:acg-uniqueness}.
Zupełność miary Lebesgue'a zapewnia mierzalność $f$:
poza zbiorem zerowym jest ona granicą funkcji mierzalnych,
a na zbiorze wyjątkowym przyjmuje dowolne skończone wartości.
Wszystkie nierówności i zbieżność zachodzą poza jednym
mierzalnym zbiorem zerowym, będącym sumą przeliczalnie wielu wyjątków. Zastępujemy na nim $f_n,f$ oraz (w przypadku z ograniczeniami) $g,h$ zerem. Nie zmienia to szerokich całek na żadnym podprzedziale ani całek Lebesgue'a funkcji już całkowalnych; po wykazaniu całkowalności granicy można przywrócić jej początkowe wartości na mocy definicji \ref{it:def:denjoy-broad} i zupełności miary Lebesgue'a.

W pierwszym przypadku $u_n=f_n-f_1$ są nieujemne
prawie wszędzie i całkowalne w sensie Lebesgue'a na mocy
twierdzenia \ref{ma1:kh:thm:wide-denjoy-positive}.
Z liniowości oraz zgodności całek mamy
\[
 \int_a^b u_n\,\mathrm d\lambda_1
 =(\mathrm D_{\mathrm{sz}})\int_a^b f_n
 -(\mathrm D_{\mathrm{sz}})\int_a^b f_1.
\]
Lewa strona jest ograniczona z góry przez założenie.
Ponieważ $u_n\uparrow u=f-f_1$ prawie wszędzie, twierdzenie
\atlasref{ME-08::thm:mct-ae}{ME-08: zbieżność monotoniczna prawie wszędzie} daje $u\in\mathcal L^1$ i
$\int |u-u_n|\,\mathrm d\lambda_1
=\int u\,\mathrm d\lambda_1-\int u_n\,\mathrm d\lambda_1\to0$.
W drugim przypadku $u_n=f_n-g$ i $w=h-g$ są nieujemne
prawie wszędzie i całkowalne w sensie Lebesgue'a przez to
samo twierdzenie o funkcjach nieujemnych. Mamy
$0\le u_n\le w$ prawie wszędzie oraz $u_n\to u=f-g$
prawie wszędzie, więc $0\le u\le w$ prawie wszędzie.
Z monotoniczności całki Lebesgue'a i $0\le u\le w$ mamy $u\in\mathcal L^1$. Twierdzenie \atlasref{ME-08::thm:dct}{ME-08: zbieżność zmajoryzowana} zastosowane do $|u_n-u|\le w$ daje $\int_{[a,b]}|u_n-u|\,\mathrm d\lambda_1\to0$.

W obu przypadkach $u$ jest całkowalna w szerokim sensie
Denjoya, ponieważ jej bezwzględnie ciągła pierwotna
Lebesgue'a należy do $\AC G$ i ma pochodną aproksymatywną
równą $u$ prawie wszędzie (twierdzenie
\ref{ma1:kh:thm:ac-lebesgue-primitive}). Dodając odpowiednio
$f_1$ lub $g$, dostajemy całkowalność $f$.
Ponadto $f_n-f=u_n-u\in\mathcal L^1$ i zachodzi
żądana zbieżność różnic w normie $L^1$. W przypadku ciągu malejącego stosujemy pierwszy przypadek do $-f_n\uparrow-f$: ograniczenie całek od dołu przechodzi w ograniczenie od góry, a $|(-f_n)-(-f)|=|f_n-f|$.
Z liniowości, lokalności i zgodności całek,
$F_n-F$ jest pierwotną Lebesgue'a funkcji $f_n-f$,
zerującą się w $a$. Wniosek \ref{it:cor:acg-bv}
daje dokładny wzór na jej wahanie.
Dla $a\le s\le t\le b$ mamy
\[
 \left|(\mathrm D_{\mathrm{sz}})\int_s^t(f_n-f)\right|
 =\left|(F_n-F)(t)-(F_n-F)(s)\right|
 \le V_a^b(F_n-F).
\]
Prawa strona nie zależy od $s,t$ i dąży do zera,
co daje wszystkie pozostałe tezy.
\end{proof}
\begin{twierdzenie}[Lemat Fatou dla szerokiej całki Denjoya]
\label{it:thm:wide-denjoy-fatou}
Niech $f_n$ oraz $g$ będą całkowalne w szerokim sensie
Denjoya na $[a,b]$ i niech $g\le f_n$ prawie wszędzie
dla każdego $n$. Załóżmy, że
\[
 L=\liminf_{n\to\infty}(\mathrm D_{\mathrm{sz}})\int_a^b f_n
 <+\infty.
\]
Wtedy $\liminf_n f_n$ jest skończona prawie wszędzie.
Jej dowolny reprezentant $f$ o skończonych wartościach,
równy tej granicy dolnej prawie wszędzie, jest całkowalny
w szerokim sensie Denjoya, $f-g\in\mathcal L^1[a,b]$ oraz
\[
 (\mathrm D_{\mathrm{sz}})\int_a^b f\le L.
\]
\end{twierdzenie}
\begin{proof}
Jeżeli $a=b$, wszystkie warunki dotyczące zachodzenia
prawie wszędzie są automatycznie spełnione, każda
skończona funkcja na $\{a\}$ jest całkowalna z całką
zero (także w sensie Lebesgue'a), a $L=0$. Teza więc
zachodzi. Załóżmy więc $a<b$.
Z twierdzenia \ref{ma1:kh:thm:wide-denjoy-positive} każda
funkcja $u_n=f_n-g$ jest całkowalna w sensie Lebesgue'a
i nieujemna prawie wszędzie. Możemy zmienić wszystkie
$u_n$ na jednym zbiorze miary zero, aby były nieujemne
wszędzie. Zgodność obu całek i liniowość dają
\[
 \int_a^b u_n\,\mathrm d\lambda_1
 =(\mathrm D_{\mathrm{sz}})\int_a^b f_n
  -(\mathrm D_{\mathrm{sz}})\int_a^b g .
\]
W szczególności $L\ge(\mathrm D_{\mathrm{sz}})\int_a^b g$,
więc $L$ jest liczbą rzeczywistą.
Nieujemna mierzalna funkcja rozszerzona $u=\liminf_n u_n$
spełnia, na mocy lematu Fatou \atlasref{ME-08::thm:fatou}{ME-08: lemat Fatou},
\[
 \int_a^b u\,\mathrm d\lambda_1
 \le \liminf_{n\to\infty}\int_a^b u_n\,\mathrm d\lambda_1
 =L-(\mathrm D_{\mathrm{sz}})\int_a^b g<+\infty.
\]
Zatem $u$ jest skończona prawie wszędzie: gdyby zbiór $\{u=+\infty\}$ miał dodatnią miarę, monotoniczność całki zastosowana do $k\chi_{\{u=+\infty\}}$ dla wszystkich $k\ge1$ dawałaby całkę nieskończoną. Po zmianie
wartości na zbiorze wyjątkowym należy do $\mathcal L^1$.
Poza zbiorem miary zero zachodzi
$\liminf_n f_n=g+u$, gdyż $g$ ma skończone wartości.
Każdy reprezentant $f$ opisany w tezie spełnia zatem
$f-g=u$ prawie wszędzie i jest mierzalny dzięki zupełności
miary Lebesgue'a. Funkcja $f-g$ jest całkowalna w sensie
Lebesgue'a, a więc również w szerokim sensie Denjoya.
Dodając $g$ i stosując liniowość oraz powyższą nierówność,
otrzymujemy
\[
 (\mathrm D_{\mathrm{sz}})\int_a^b f
 =(\mathrm D_{\mathrm{sz}})\int_a^b g
  +\int_a^b (f-g)\,\mathrm d\lambda_1\le L.
\]
\end{proof}
\begin{przyklad}[Konieczność ograniczenia od dołu w lemacie Fatou]
Na $[0,1]$ połóżmy $f_n=-n\chi_{(0,1/n)}$ dla $n\ge1$.
Każda z tych funkcji jest całkowalna w sensie Lebesgue'a
i ma całkę $-1$, również w szerokim sensie Denjoya.
Dla każdego $x\in[0,1]$ mamy $f_n(x)\to0$: w punkcie
$0$ wszystkie wartości są zerowe, a dla $x>0$ od
pewnego indeksu $1/n<x$. Całka granicy dolnej wynosi
zatem $0>-1$.
Nie istnieje wspólne ograniczenie od dołu $g$ całkowalne
w szerokim sensie Denjoya. Istotnie, gdyby istniało,
to $g\le f_1\le0$ prawie wszędzie, więc z twierdzenia
\ref{ma1:kh:thm:wide-denjoy-positive} zastosowanego do
$-g$ mielibyśmy $-g\in\mathcal L^1$.
Poza jednym zbiorem miary zero na każdym przedziale
$(1/(k+1),1/k)$ dla $k\ge1$ zachodziłoby $-g\ge k$.
Przeliczalna addytywność całki nieujemnej dawałaby
\[
 \int_0^1(-g)\,\mathrm d\lambda_1
 \ge \sum_{k\ge1} k\left(\frac1k-\frac1{k+1}\right)
 =\sum_{k\ge1}\frac1{k+1}=+\infty,
\]
co jest sprzecznością.
\end{przyklad}

\begin{twierdzenie}[Podstawienie w szerokiej całce Denjoya]\label{ma1:kh:thm:wide-denjoy-substitution}
Niech $a<b$, $u<v$, a $\phi:[u,v]\to[a,b]$ będzie dyfeomorfizmem klasy $C^1$ z pochodną nigdzie niezerową. Jeżeli $f$ jest całkowalna w szerokim sensie Denjoya na $[a,b]$, to $(f\circ\phi)\phi'$ jest całkowalna w tym sensie na $[u,v]$ i
\[
 (\mathrm D_{\mathrm{sz}})\int_u^v f(\phi(t))\phi'(t)\,\mathrm dt
 =(\mathrm D_{\mathrm{sz}})\int_{\phi(u)}^{\phi(v)}f(x)\,\mathrm dx.
\]
Całkę o odwróconych granicach rozumiemy ze znakiem minus.
\end{twierdzenie}
\begin{proof}
Niech $F$ będzie pierwotną $f$ klasy $\AC G$ i niech $E_n$ będzie jej domkniętym pokryciem. Połóżmy $M=\max|\phi'|$, $m=\min|\phi'|>0$. Z twierdzenia Lagrange'a
$m|s-t|\le|\phi(s)-\phi(t)|\le M|s-t|$.
Zbiory $\phi^{-1}(E_n)$ są domknięte i pokrywają $[u,v]$. Obrazy przedziałów o rozłącznych wnętrzach mają rozłączne wnętrza, końce w $E_n$ i łączną długość nie większą niż $M$ razy suma długości przedziałów wyjściowych. Moduły przyrostów $F\circ\phi$ są dokładnie modułami przyrostów $F$ na obrazach (także gdy $\phi$ jest malejąca). Dokładniej, dla ustalonego $\varepsilon>0$ wybieramy
$\delta>0$ z warunku $\AC(E_n)$ dla błędu $\varepsilon$.
Jeżeli suma długości przedziałów wyjściowych jest
mniejsza od $\delta/M$, to suma długości ich obrazów
jest mniejsza od $\delta$. Po uporządkowaniu końców
obrazów rosnąco warunek $\AC(E_n)$ ogranicza sumę
modułów przyrostów przez $\varepsilon$.
Stąd $F\circ\phi\in\AC(\phi^{-1}(E_n))$ dla każdego $n$;
ciągłość złożenia oraz domknięte pokrycie dają
$F\circ\phi\in\AC G$.

Ustalmy $t\in(u,v)$, dla którego istnieje $F'_{\mathrm{ap}}(\phi(t))=L$. Dla każdego $\eta>0$ zbiór złych ilorazów $B_\eta$ funkcji $F$ z progiem $\eta$ ma gęstość zero w $x=\phi(t)$. Dla małego $r>0$ obraz $(t-r,t+r)$ zawiera się w $(x-Mr,x+Mr)$. Odwrotność $\phi^{-1}$ jest lipschitzowska ze stałą $1/m$.
Dla dowolnego $A\subset[a,b]$ i pokrycia $A$ przedziałami
$J_k$ zbiory $\phi^{-1}(J_k\cap[a,b])$ są przedziałami
(o ile są niepuste) o długościach nie większych od $|J_k|/m$.
Pokrywają $\phi^{-1}(A)$, więc po wzięciu infimum po pokryciach
otrzymujemy $\lambda_1^*(\phi^{-1}(A))\le\lambda_1^*(A)/m$. Zbiór $B_\eta$ jest borelowski na mocy ciągłości $F$, a jego przeciwobraz przez ciągłą $\phi$ również jest borelowski. Stąd
\[
 \frac{\lambda_1(\phi^{-1}(B_\eta)\cap(t-r,t+r))}{2r}
 \le \frac{M}{m}
 \frac{\lambda_1(B_\eta\cap(x-Mr,x+Mr))}{2Mr}\longrightarrow0.
\]
Poza tym zbiorem iloraz różnicowy $F\circ\phi$ różni się od $L\phi'(t)$ o co najwyżej $M\eta+|L|\,|(\phi(s)-\phi(t))/(s-t)-\phi'(t)|$. Dla zadanego $\varepsilon>0$ wybieramy najpierw
$\eta=\varepsilon/(2(1+M))$, a następnie promień
otoczenia $t$ tak mały, aby
\[
 \left|\frac{\phi(s)-\phi(t)}{s-t}-\phi'(t)\right|
 <\frac{\varepsilon}{2(1+|L|)}
 \quad (s\ne t).
\]
W tym otoczeniu błąd jest mniejszy od $\varepsilon$
poza $\phi^{-1}(B_\eta)$, który ma gęstość zero w $t$.
Zbiór złych ilorazów dla progu $\varepsilon$ ma zatem
gęstość zero w $t$. To dowodzi reguły łańcuchowej
dla pochodnej aproksymatywnej w rozważanych punktach. Przeciwobraz zbioru zerowego, na którym $F'_{\mathrm{ap}}=f$ może nie zachodzić, jest zerowy przez to samo oszacowanie lipschitzowskie. Zatem $(F\circ\phi)'_{\mathrm{ap}}=(f\circ\phi)\phi'$ prawie wszędzie. Twierdzenie \ref{tw:denjoy-szeroka-ftc} daje żądaną równość przyrostów.
\end{proof}

\begin{twierdzenie}[Mnożniki o ograniczonym wahaniu w szerokiej całce Denjoya]\label{ma1:kh:thm:wide-denjoy-multiplier}
Niech $f$ będzie całkowalna w szerokim sensie Denjoya, $F\in\AC G[a,b]$ jej pierwotną, a $g\in\BV[a,b]$. Wówczas $fg$ jest całkowalna w szerokim sensie Denjoya oraz
\[
 (\mathrm D_{\mathrm{sz}})\int_a^b fg
 =F(b)g(b)-F(a)g(a)-\int_a^b F\,\mathrm dg,
\]
gdzie ostatnia całka jest całką Riemanna--Stieltjesa.
\end{twierdzenie}
\begin{proof}
Ponieważ $F$ jest ciągła, całki Riemanna--Stieltjesa z $F$ względem $g$ istnieją na wszystkich podprzedziałach (twierdzenie \ref{it:thm:rs-tools}). Definiujemy
$H(x)=F(x)g(x)-\int_a^x F\,\mathrm dg$.
Z addytywności i liniowości całki Riemanna--Stieltjesa oraz całkowania funkcji stałej, dla $u<v$ mamy
\[
 H(v)-H(u)=g(v)(F(v)-F(u))-\int_u^v(F(z)-F(u))\,\mathrm dg(z).
\]
Oznaczmy moduł ciągłości $F$ przez $\omega_F(s)=\sup\{|F(x)-F(y)|:x,y\in[a,b],\ |x-y|\le s\}$. Z jednostajnej ciągłości $F$ mamy $\omega_F(s)\to0$ przy $s\downarrow0$. Stąd
\[
 |H(v)-H(u)|\le \|g\|_\infty |F(v)-F(u)|
   +\omega_F(v-u)\Var_{[u,v]}g.
\]
Dowodzi to ciągłości $H$ także w punktach skoków $g$. Jeśli $F\in\AC(E_n)$, dla rodziny przedziałów o końcach w $E_n$, rozłącznych wnętrzach i łącznej długości mniejszej od $\eta$ suma modułów przyrostów $H$ nie przekracza sumy modułów przyrostów $F$ pomnożonej przez $\|g\|_\infty$ oraz $\omega_F(\eta)\Var_{[a,b]}g$. W drugim składniku używamy addytywności wahania (lemat \ref{ma1:kh:lem:variation-additivity}). Dla $\varepsilon>0$ wybieramy $\eta$ z warunku $\AC(E_n)$ dla błędu $\varepsilon/(2(1+\|g\|_\infty))$ i zmniejszamy je tak, by $\omega_F(\eta)\Var_{[a,b]}g<\varepsilon/2$. Zatem $H\in\AC G[a,b]$.

Niech $V(x)=\Var_{[a,x]}g$. Funkcja $V$ jest
niemalejąca i ograniczona. Uzasadnimy, że ma skończoną
pochodną prawie wszędzie, również gdy nie jest
prawostronnie ciągła. Przedłużamy $V$ przez stałe
poza $[a,b]$ i oznaczamy $W(t)=V(t+)$.
Monotoniczność daje istnienie tych granic, a $W$
jest niemalejąca i prawostronnie ciągła: istotnie,
$W(t)=\inf_{s>t}V(s)$, a dla $t<r<s$ mamy
$W(t)\le W(r)\le V(s)$. Przejście z $r\downarrow t$
i następnie infimum po $s>t$ daje prawostronną
ciągłość. Z twierdzenia \atlasref{ME-04::it:thm:stieltjes-construction}{ME-04: konstrukcja miary Lebesgue'a--Stieltjesa}
funkcja $W$ wyznacza skończoną dodatnią miarę
Lebesgue'a--Stieltjesa, gdyż jest stała poza przedziałem
zwartym. Jej funkcja rozkładu różni się od $W$
o stałą $\lim_{t\to-\infty}W(t)$, więc z twierdzenia
\atlasref{ME-25::thm:cdf-derivative-25}{ME-25: pochodna funkcji rozkładu} pochodna $W'$ istnieje
i jest skończona prawie wszędzie.

Funkcja $V$ ma co najwyżej przeliczalnie wiele
punktów nieciągłości, jak każda funkcja niemalejąca.
Ustalmy punkt $x\in(a,b)$ ciągłości $V$, w którym
istnieje skończona pochodna $W'(x)=D$.
Wówczas $W(x)=V(x)$, a dla $h\ne0$ dostatecznie
bliskiego zeru monotoniczność daje
\[
 W(x+h-h^2)\le V(x+h)\le W(x+h).
\]
Oba skrajne wyrażenia mają postać $V(x)+Dh+o(|h|)$,
ponieważ $(h-h^2)/h\to1$ i $W$ jest różniczkowalna w $x$.
Zatem $V(x+h)=V(x)+Dh+o(|h|)$, czyli $V'(x)=D$.
To dowodzi istnienia skończonej pochodnej $V$
prawie wszędzie bez dodatkowego warunku ciągłości. Funkcja $g$ jest ciągła poza zbiorem co najwyżej przeliczalnym, co wynika z rozkładu Jordana na różnicę funkcji niemalejących. W punkcie $x$, w którym $g$ jest ciągła, $V'$ istnieje i jest skończona oraz $F'_{\mathrm{ap}}(x)=f(x)$, dla $y\ne x$ (z konwencją całki o odwróconych granicach) mamy
\[
 \frac{H(y)-H(x)}{y-x}
 =g(y)\frac{F(y)-F(x)}{y-x}
 -\frac{1}{y-x}\int_x^y(F(z)-F(x))\,\mathrm dg(z).
\]
Moduł ostatniego składnika jest nie większy niż $\omega_F(|y-x|)|V(y)-V(x)|/|y-x|$, co dąży do zera. Pierwszy składnik ma granicę aproksymatywną $g(x)f(x)$. Istotnie, pisząc $Q(y)=(F(y)-F(x))/(y-x)$, mamy
\[
 |g(y)Q(y)-g(x)f(x)|
 \le \|g\|_\infty|Q(y)-f(x)|+|f(x)|\,|g(y)-g(x)|.
\]
Dla $\varepsilon>0$ pierwszy składnik jest mniejszy od $\varepsilon/2$ poza zbiorem gęstości zero, z definicji pochodnej aproksymatywnej zastosowanej z progiem $\varepsilon/(2(1+\|g\|_\infty))$. Drugi jest mniejszy od $\varepsilon/2$ w dostatecznie małym otoczeniu $x$, z ciągłości $g$ w $x$. Zbiór, w którym łączny błąd jest co najmniej $\varepsilon$, ma więc gęstość zero w $x$. Zatem $H'_{\mathrm{ap}}=fg$ prawie wszędzie. Wzór wynika z twierdzenia \ref{tw:denjoy-szeroka-ftc}.
\end{proof}

\begin{wniosek}[Całkowanie przez części z mnożnikiem bezwzględnie ciągłym]\label{ma1:kh:cor:wide-denjoy-ibp-ac}
W hipotezach twierdzenia \ref{ma1:kh:thm:wide-denjoy-multiplier}, jeśli $g\in\AC[a,b]$, to
\[
 (\mathrm D_{\mathrm{sz}})\int_a^b fg
 =F(b)g(b)-F(a)g(a)-\int_a^b F(x)g'(x)\,\mathrm d\lambda_1(x).
\]
\end{wniosek}
\begin{proof}
Bezwzględna ciągłość daje ograniczone wahanie i reprezentację $g(v)-g(u)=\int_u^v g'\,\mathrm d\lambda_1$ z $g'\in L^1$ (twierdzenie \ref{ma1:kh:thm:ac-lebesgue-primitive}). Dla podziału znaczonego $P=\{(t_j,[x_{j-1},x_j])\}$
reprezentacja przyrostów $g$ daje
\[
\begin{aligned}
 S(F,g;P)-\int_{[a,b]}F(x)g'(x)\,\mathrm d\lambda_1(x)
 &=\sum_j\int_{[x_{j-1},x_j]}
 (F(t_j)-F(x))g'(x)\,\mathrm d\lambda_1(x).
\end{aligned}
\]
Końce komórek mają miarę zero, a $Fg'$ jest całkowalna
Lebesgue'a, ponieważ $F$ jest ograniczona i $g'\in L^1$.
Moduł prawej strony nie przekracza
$\omega_F(\mesh P)\int_{[a,b]}|g'|\,\mathrm d\lambda_1$,
co dąży do zera przy $\mesh P\to0$.
Z definicji całki Riemanna--Stieltjesa otrzymujemy
$\int_a^b F\,\mathrm dg=\int_{[a,b]}Fg'\,\mathrm d\lambda_1$. Stosujemy twierdzenie \ref{ma1:kh:thm:wide-denjoy-multiplier}.
\end{proof}

\begin{twierdzenie}[Dołączenie końca przedziału w szerokiej całce Denjoya]\label{ma1:kh:thm:wide-denjoy-hake}
Niech $a<b$, a $f:[a,b]\to\R$ ma skończone wartości. Całka w szerokim sensie Denjoya na $[a,b]$ istnieje wtedy i tylko wtedy, gdy istnieje na każdym $[c,b]$, $a<c<b$, oraz istnieje skończona granica
$L=\lim_{c\downarrow a}(\mathrm D_{\mathrm{sz}})\int_c^b f$.
Wtedy całka na $[a,b]$ jest równa $L$.
\end{twierdzenie}
\begin{proof}
Konieczność wynika z ciągłości pierwotnej $G$ i wzoru \ref{tw:denjoy-szeroka-ftc}: całka po $[c,b]$ jest równa $G(b)-G(c)$, co dąży do $G(b)-G(a)$. Odwrotnie, definiujemy $F(a)=0$ i $F(x)=L-(\mathrm D_{\mathrm{sz}})\int_x^b f$ dla $x>a$. Na każdym $[c,b]$, jeżeli $G_c$ jest pierwotną $f$ na tym podprzedziale, to z twierdzenia \ref{tw:denjoy-szeroka-ftc} mamy $F(x)=G_c(x)+L-G_c(b)$ dla $x\in[c,b]$. Jest tam więc ciągła i klasy $\AC G$. Ciągłość w $a$ wynika z definicji $L$. Wybieramy ciąg $c_n\downarrow a$ z $a<c_n<b$. Każdy $[c_n,b]$ ma przeliczalne pokrycie przez zbiory domknięte, na których $F$ spełnia $\AC$; są one również domknięte w $[a,b]$. Łącząc te pokrycia ze zbiorem $\{a\}$, otrzymujemy przeliczalne domknięte pokrycie całego $[a,b]$, a więc $F\in\AC G[a,b]$. Poza przeliczalną sumą lokalnych zbiorów wyjątkowych oraz zbioru $\{a,b,c_1,c_2,\ldots\}$ mamy $F'_{\mathrm{ap}}=f$, ponieważ każdy pozostały punkt należy do wnętrza któregoś $[c_n,b]$ i pochodna aproksymatywna zależy tylko od wartości w jego otoczeniu. Stąd całka istnieje z wartością $F(b)-F(a)=L$.
\end{proof}

\begin{wniosek}[Dołączenie prawego końca i obu końców przedziału]
\label{it:cor:wide-denjoy-hake-two-endpoints}
Niech $a<b$ i $f:[a,b]\to\R$ ma skończone wartości.
Całka w szerokim sensie Denjoya na $[a,b]$ istnieje wtedy
i tylko wtedy, gdy istnieje na każdym $[a,d]$, $a<d<b$,
oraz istnieje skończona granica
$\lim_{d\uparrow b}(\mathrm D_{\mathrm{sz}})\int_a^d f$.
Całka na $[a,b]$ jest wówczas równa tej granicy.
Jeżeli $f$ jest całkowalna na każdym zwartym podprzedziale
$(a,b)$ i $s\in(a,b)$, to całkowalność na $[a,b]$
jest równoważna istnieniu obu skończonych granic
\[
 L_-=\lim_{c\downarrow a}(\mathrm D_{\mathrm{sz}})\int_c^s f,
 \qquad
 L_+=\lim_{d\uparrow b}(\mathrm D_{\mathrm{sz}})\int_s^d f.
\]
Wtedy $(\mathrm D_{\mathrm{sz}})\int_a^b f=L_-+L_+$.
\end{wniosek}
\begin{proof}
Dla prawego końca stosujemy twierdzenie
\ref{ma1:kh:thm:wide-denjoy-hake} do $h(t)=f(a+b-t)$.
Z twierdzenia o podstawieniu
\ref{ma1:kh:thm:wide-denjoy-substitution} dla odbicia
$\phi(t)=a+b-t$, o pochodnej $-1$, i z liniowości
wynika, że dla $a<c<b$ całkowalność $h$ na $[c,b]$
jest równoważna całkowalności $f$ na $[a,a+b-c]$,
a ich całki są równe. Kierunek odwrotny wynika
z ponownego zastosowania tego samego odbicia.
Analogicznie dla całego $[a,b]$.
Gdy $c\downarrow a$, mamy $a+b-c\uparrow b$,
co dowodzi pierwszej części.

Dla dwóch końców konieczność wynika z lokalności
i ciągłości pierwotnej $F$: granice są odpowiednio
równe $F(s)-F(a)$ oraz $F(b)-F(s)$.
Odwrotnie, twierdzenie \ref{ma1:kh:thm:wide-denjoy-hake}
na $[a,s]$ daje całkowalność z wartością $L_-$,
a pierwsza część niniejszego wniosku na $[s,b]$
daje całkowalność z wartością $L_+$.
Stwierdzenie \ref{st:denjoy-szeroka-lokalnosc}
daje całkowalność na całości i sumę obu wartości.
\end{proof}
\begin{uwaga}[Symetryczne obcięcie nie wystarcza]
Niech $f(x)=1/x-1/(1-x)$ dla $0<x<1$ i $f(0)=f(1)=0$.
Funkcja jest ciągła na $(0,1)$, a jej pierwotna na
tym przedziale to $\log x+\log(1-x)$.
Dla $0<\varepsilon<1/2$ mamy
\[
 \int_\varepsilon^{1-\varepsilon} f(x)\,\mathrm dx=0,
 \qquad
 \int_\varepsilon^{1/2}f(x)\,\mathrm dx
 =\log(1/4)-\log(\varepsilon(1-\varepsilon))
 \longrightarrow+\infty.
\]
Całka po $[1/2,1-\varepsilon]$ jest przeciwna do
drugiej z tych całek, więc dąży do $-\infty$.
Na tych zwartych podprzedziałach całki Riemanna
i szeroka Denjoya są zgodne.
Mimo istnienia granicy przy symetrycznym obcięciu
całka szeroka na $[0,1]$ nie istnieje, ponieważ
dwie osobne granice z poprzedniego wniosku nie są skończone.
\end{uwaga}

\begin{lemat}[Konieczny warunek na osobliwość końcową]\label{it:lem:broad-endpoint-necessary}
Niech $a<b$ i $f:[a,b]\to\R$ będzie ciągła na $(a,b]$. Jeżeli jest całkowalna w szerokim sensie Denjoya na $[a,b]$, to istnieje skończona granica
\[
 \lim_{c\downarrow a}\int_c^b f(x)\,\mathrm dx,
\]
gdzie dla $c>a$ występuje całka Riemanna.
\end{lemat}
\begin{proof}
Niech $F\in\AC G[a,b]$ będzie jej ciągłą pierwotną. Z twierdzenia \ref{tw:denjoy-szeroka-ftc} oraz zgodności szerokiej całki Denjoya z KH i całką Riemanna otrzymujemy dla $a<c<b$ równość $\int_c^b f(x)\,\mathrm dx=F(b)-F(c)$. Ciągłość $F$ w $a$ daje skończoną granicę $F(b)-F(a)$. Wniosek dotyczy także KH i wąskiej całki Denjoya, ponieważ są zawarte w szerokiej całce Denjoya.
\end{proof}
\subsection{Przykłady parametryczne i kryteria całkowalności}
\begin{przyklad}[Dodatnia osobliwość potęgowa]\label{ma1:kh:ex:positive-power}
Niech $f_p(x)=x^{-p}$ dla $x\in(0,1]$ oraz $f_p(0)=0$. Na $[0,1]$ funkcja jest całkowalna w sensie Riemanna dokładnie dla $p\le0$, natomiast w sensie Lebesgue'a, McShane'a, KH i Denjoya dla $p<1$. Istotnie,
\[
\int_c^1x^{-p}\,\mathrm dx=\begin{cases}(1-c^{1-p})/(1-p),&p\ne1,\\-\log c,&p=1.\end{cases}
\]
Granica przy $c\downarrow0$ jest skończona dokładnie dla $p<1$, a z twierdzenia \atlasref{ME-08::thm:mct}{ME-08: twierdzenie Beppo Leviego} zastosowanego do $f_p\chi_{[1/n,1]}$ wynika ten sam warunek całkowalności Lebesgue'a. Rachunek całki potęgowej jest szczególnym przypadkiem przykładu \atlasref{ME-08::ex:power-integral}{ME-08: całka potęgi blisko zera}. Dla $p<1$ stosujemy zgodność Lebesgue'a z McShane'em, KH i obiema całkami Denjoya; dla $p\ge1$ lemat \ref{it:lem:broad-endpoint-necessary} wyklucza również szeroką całkę Denjoya. Gdy $p>0$, funkcja jest nieograniczona; gdy $p\le0$, jest ograniczona i ma co najwyżej jeden punkt nieciągłości, więc stosujemy twierdzenie \atlasref{ME-08::thm:riemann-criterion}{ME-08: kryterium Lebesgue'a}.
\end{przyklad}

\begin{przyklad}[Oscylacja zmienia próg]\label{ma1:kh:ex:oscillatory-thresholds}
Niech $g_p(x)=x^{-p}\sin(1/x)$ dla $x\in(0,1]$ i $g_p(0)=0$.
Wtedy
\[
\begin{array}{c|c}
\text{Teoria całkowania}&\text{Warunek całkowalności na }[0,1]\\\hline
\text{Riemann / Darboux}&p\le0\\
\text{Lebesgue / McShane}&p<1\\
\text{KH / Denjoy (obie wersje)}&p<2
\end{array}
\]
Twierdzenie \ref{ma1:kh:tw:hk-oscylacyjne} dla $q=1$
daje progi $p<1$ i $p<2$ dla całek Lebesgue'a i KH.
W tym samym dowodzie wykazano rozbieżność całki niewłaściwej
przy $p\ge2$; lemat \ref{it:lem:broad-endpoint-necessary}
wyklucza wtedy również całkowalność w szerokim sensie Denjoya.
Dla $p<2$ stosujemy inkluzję KH w obie klasy Denjoya,
a dla McShane'a równoważność z całkowalnością Lebesgue'a.
Jeśli $p\le0$, funkcja jest ograniczona i ciągła poza
co najwyżej jednym punktem, więc stosujemy kryterium
\atlasref{ME-08::thm:riemann-criterion}{ME-08: kryterium Lebesgue'a}.
Dla $p>0$ wartości w punktach
$x_k=(\pi/2+2k\pi)^{-1}$ dążą do nieskończoności.
Brak ograniczoności wyklucza całkę Riemanna i Darboux.
\end{przyklad}

\begin{przyklad}[Osobliwość logarytmiczna]\label{ma1:kh:ex:log-singularity}
Dla $q\in\R$ definiujemy $h_q(x)=1/[x(\log(1/x))^q]$ dla $0<x\le e^{-1}$ i $h_q(0)=0$. Funkcja jest nieograniczona przy $0$ dla każdego $q$, ponieważ po podstawieniu $u=\log(1/x)$ jej wartości to $e^u/u^q\to+\infty$, więc nie jest całkowalna w sensie właściwego Riemanna. Jest całkowalna w sensie Lebesgue'a, McShane'a, KH i Denjoya wtedy i tylko wtedy, gdy $q>1$. Po podstawieniu $u=\log(1/x)$ otrzymujemy
$\int_c^{e^{-1}}h_q(x)\,\mathrm dx=\int_1^{\log(1/c)}u^{-q}\,\mathrm du$, która ma granicę skończoną dokładnie dla $q>1$ (dla $q=1$ jest to logarytm, a dla $q\ne1$ stosujemy wzór potęgowy). Dla $q>1$ z twierdzenia \atlasref{ME-08::thm:mct}{ME-08: twierdzenie Beppo Leviego} i zgodności całek otrzymujemy wszystkie wymienione całki, a dla $q\le1$ lemat \ref{it:lem:broad-endpoint-necessary} wyklucza także szeroką całkę Denjoya.
\end{przyklad}

\begin{przyklad}[Ograniczona funkcja niecałkowalna Riemanna]\label{ma1:kh:ex:fat-cantor-no-riemann}
Niech $K\subset[0,1]$ będzie domkniętym zbiorem o dodatniej mierze Lebesgue'a i pustym wnętrzu (grubym zbiorem Cantora). Taki zbiór otrzymamy, numerując wymierne punkty $[0,1]$ jako $(q_j)_{j\ge1}$ i usuwając z $[0,1]$ sumę względnie otwartych przedziałów wokół $q_j$ o długościach co najwyżej $2^{-j-2}$. Pozostały zbiór jest domknięty, nie zawiera żadnego względnie otwartego niepustego przedziału i ma miarę co najmniej $1-\sum_{j\ge1}2^{-j-2}=3/4$. Wówczas $\chi_K$ jest ograniczona i całkowalna w sensie Lebesgue'a, z całką $\lambda_1(K)>0$, lecz na mocy twierdzenia \atlasref{ME-08::thm:riemann-criterion}{ME-08: kryterium Lebesgue'a} nie jest całkowalna w sensie Riemanna: jej zbiór punktów nieciągłości to $K$, mający dodatnią miarę. Nie można usunąć tej przeszkody przez zmianę wartości jedynie na zbiorze miary zero: w prawie każdym punkcie gęstości $K$ są arbitralnie bliskie punkty z wartością $1$, a dzięki temu, że dopełnienie $K$ jest otwarte i gęste, każde otoczenie zawiera podprzedział dodatniej długości w dopełnieniu, a więc również punkty z wartością $0$ poza zbiorem zmienionych wartości. Dokładniej, jeśli $N$ jest zbiorem zmiany wartości, to dla każdego punktu gęstości jeden zbioru $K$ i każdego dostatecznie małego promienia $r$ mamy $\lambda_1(K\cap B(x,r))>0$. Zatem $(K\setminus N)\cap B(x,r)\ne\varnothing$. W tym samym otoczeniu istnieje również punkt $( [0,1]\setminus K)\setminus N$, jak ustalono wyżej. Zmieniona funkcja ma więc oscylację co najmniej $1$ w każdym takim punkcie (niezależnie od jej wartości w samym punkcie). Twierdzenie \atlasref{ME-25::cor:lebesgue-density-25}{ME-25: twierdzenie o punktach gęstości} daje zbiór takich punktów o mierze $\lambda_1(K)>0$, co wyklucza całkowalność Riemanna.
\end{przyklad}

\begin{uwaga}[Nazewnictwo i oddzielne poziomy teorii]
W terminologii historycznej szeroka całka Denjoya jest także nazywana całką Denjoya--Chinczyna. Jej opis wymaga $\AC G$ i pochodnej aproksymatywnej. Wąska całka Denjoya, równoważna KH i Perronowi, wymaga $\AC G_*$ oraz zwykłej pochodnej prawie wszędzie. Jeżeli $F\in\AC G[a,b]$ ma zwykłą skończoną pochodną prawie wszędzie, to każda funkcja o skończonych wartościach równa tej pochodnej prawie wszędzie jest całkowalna w szerokim sensie Denjoya, gdyż zwykła pochodna jest również aproksymatywna. Sama ciągłość i istnienie zwykłej pochodnej prawie wszędzie nie wystarczają (funkcja Cantora). Nie wynika też stąd równoważność klas $\AC G$ i $\AC G_*$; przykład \ref{ma1:kh:ex:denjoy-strict} pokazuje już, że nawet przy zwykłej pochodnej prawie wszędzie klasa $\AC G$ nie gwarantuje całkowalności KH.
\end{uwaga}
\endgroup
